\subsection{28.04.2026}

\subsubsection{Ziel}

Ziel des Tages war es, die Taupunktberechnung
funktionsfähig umzusetzen.

Zusätzlich sollte eine Datenbank vorbereitet werden,
um die Messwerte langfristig zu speichern.
Außerdem sollte recherchiert werden, wie der Lüfter
angeschlossen und angesteuert werden kann.

\subsubsection{Durchgeführte Arbeiten}

Zunächst wurden mehrere Probleme innerhalb der
Taupunktberechnung identifiziert, welche zu
Fehlern während der Berechnung führten.

Zur Fehlersuche wurden verschiedene
Zwischenergebnisse auf der Konsole ausgegeben,
um die genaue Fehlerquelle besser eingrenzen zu können.

Zusätzlich wurde die Ausgabe bei ungültigen
Messwerten angepasst.
Die Meldung „Messung läuft...“ wird nun nur noch
alle 20 ungültigen Messwerte ausgegeben.
Nach 100 ungültigen Messversuchen wird der
Vorgang mit der Meldung
„Messung war nicht möglich“
abgebrochen.

Darüber hinaus wurden einzelne Zwischenergebnisse
der Taupunktberechnung validiert.
Falls ein mathematisch ungültiger Wert entsteht,
wird nun eine Fehlermeldung ausgegeben und
der Wert 0 zurückgegeben.

Ein Teil der Fehler innerhalb der
Taupunktberechnung konnte dadurch bereits behoben
werden.

Zusätzlich wurde SQLite als mögliche
Datenbanklösung ausgewählt und erste Tests zur
Verwendung durchgeführt.

Des Weiteren wurde recherchiert, wie der Lüfter
angeschlossen und über den Raspberry Pi gesteuert
werden kann.

Nach dem Unterricht wurde zusätzlich weiter an
einer Lösung für den Anschluss des Lüfters gearbeitet.

\subsubsection{Problem 1: Fehler in der Taupunktberechnung}

Zunächst trat weiterhin ein Fehler bei der
Berechnung mit \texttt{log10()} auf.

Zusätzlich wurden Probleme beim Zusammensetzen
von Zeichenketten für die Konsolenausgabe
festgestellt.

\textbf{Lösung}

Die Konstante \texttt{c} wurde von einer
Klassenvariable zu einer Instanzvariable der
Klasse \texttt{TaupunktBerechnung} geändert.

Außerdem wurde festgestellt, dass numerische
Werte vor der Ausgabe in der Konsole zunächst
mit \texttt{str()} in Zeichenketten umgewandelt
werden müssen.

Darüber hinaus wurde erkannt, dass einige
DHT11-Sensorobjekte mehrfach und unnötig in
\texttt{Taupunkt\_main.py} erstellt wurden.

Zusätzlich wurde versehentlich das ursprüngliche
Modul \texttt{dht11} anstelle der eigenen
Implementierung aus \texttt{dht11\_jh}
importiert.

Diese Fehler wurden korrigiert.

\subsubsection{Problem 2: Ungewöhnliche Messwerte}

Die gemessenen Temperaturwerte erschienen
unerwartet niedrig.

\textbf{Lösung}

An diesem Tag konnte noch keine eindeutige
Ursache für die fehlerhaften Messwerte gefunden
werden.

\subsubsection{Problem 3: Dynamische Datenbank-Statements}

Es wurde versucht, die SQL-Statements dynamischer
aufzubauen, um doppelte Codeabschnitte zu vermeiden.

\textbf{Lösung}

Für dieses Problem konnte an diesem Tag noch
keine endgültige Lösung gefunden werden.

\subsubsection{Erkenntnisse}

Es wurde deutlich, dass mathematische
Berechnungen umfangreiche Validierungen der
Zwischenergebnisse benötigen, um Fehler frühzeitig
abzufangen.

Außerdem wurde erkannt, dass SQLite eine
geeignete Möglichkeit darstellt, Sensordaten
dauerhaft lokal auf dem Raspberry Pi zu speichern.

\subsubsection{Nächste Schritte}

Als nächster Schritt soll untersucht werden,
wie SQL-Statements mit Parametern dynamischer
gestaltet werden können.

Außerdem soll der Lüfter korrekt angeschlossen,
angesteuert und in die bestehende Hauptlogik
integriert werden.